\homework{5}{Tue 03 Oct 2023 22:01}{}

\begin{problem}
	
Consider a function, known as the correlation function, defined by the expectation value
\[
C(t)=\langle X(t) X(0)\rangle
\]
where $X(t)$ is the position operator in the Heisenberg picture. Evaluate the correlation function explicitly for the ground state of a one-dimensional harmonic oscillator.
\end{problem}
\begin{proof}[Solution]
	Recall the following identities:
	\[
		X(t) = \mathcal{U}(t)^\dag x \mathcal{U}(t)
	.\] 
	\[
		X(0) = x
	.\] 
	\[
		\mathcal{U}(t) = \exp\left( \frac{-i H t}{\hbar} \right) 
	.\] 
	\[
		x = \sqrt{\frac{\hbar}{2 m \omega}}  (a + a^{\dag} )
	.\] 
	\[
		p = i \sqrt{\frac{m \hbar \omega}{2}} (- a + a^{\dag} )
	.\] 
	\[
		N = a^{\dag}  a \qquad a a^{\dag}  = N + 1
	.\] 
	The expectation value in Heiseberg picture is
	\[
		C(t) = \left\langle 0 \right| X(t) X(0) \left| 0 \right\rangle 
	.\] 
	By eqn. (2.171),
	\[
		\exp\left( \frac{i H t}{\hbar} \right) x \exp\left( \frac{- i H t}{\hbar} \right) 
		= x \cos \omega t + \frac{p}{m \omega} \sin \omega t
	.\] 
	Use the identities to express $X(t) $ and $X(0)$ in terms of $a$ and $a^{\dag} $:
	\begin{align*}
		C(t) &= \left\langle X(t) X(0) \right\rangle  \\
		&= \left\langle 0 \right| \exp\left( \frac{i H t}{\hbar} \right) x \exp\left( \frac{- i H t}{\hbar} \right) x \left| 0 \right\rangle  \\
		&= \left\langle 0 \right| \left( x \cos \omega t + \frac{p}{m \omega} \sin \omega t \right) x \left| 0 \right\rangle  \\
		&= \cos \omega t \left\langle 0 \right| x^2 \left| 0 \right\rangle + \frac{\sin \omega t}{m \omega} \left\langle 0 \right| p x \left| 0 \right\rangle 
	.\end{align*}
	Now note that $x^2 = \frac{\hbar}{2 m \omega}\left(a^2 + \left( a^{\dag}  \right) ^2 + 2 N + 1\right)$ and $px = i \frac{\hbar}{2}\left( -a^2 + (a^{\dag} )^2 - 1 \right) $. All terms involving only $a$ or only $a^{\dag} $ do not contribute to the expectation. So
	\begin{align*}
		C(t)
		&= \cos \omega t \frac{\hbar}{2 m \omega} \left\langle 0 \right| (2N+1) \left| 0 \right\rangle + \frac{\sin \omega t}{m \omega} i \frac{\hbar}{2} \left\langle 0 \right| (-1) \left| 0 \right\rangle  \\
		&= \cos \omega t \frac{\hbar}{2 m \omega} 
		- i\frac{\sin \omega t}{m \omega} \frac{\hbar}{2}
	.\end{align*}
\end{proof}


\begin{problem}
	
For the one-dimensional simple harmonic oscillator, calculate the expectation values of $X, P, X^2$, and $P^2$ for an arbitrary state $|n\rangle$ and show that
\[
\left\langle(\Delta X)^2\right\rangle\left\langle(\Delta P)^2\right\rangle=\left(n+\frac{1}{2}\right)^2 \hbar^2
\]
\end{problem}
\begin{proof}
	First or all $\left\langle x \right\rangle  = \left\langle p \right\rangle = 0$ since both of them consists of annihilators and generators that does not preserve any eigenstate of $N$. 
	Therefore $\left\langle (\Delta x)^2 \right\rangle  = \left\langle x^2 \right\rangle  $ 
	and $\left\langle \left( \Delta p \right) ^2 \right\rangle  = \left\langle p^2 \right\rangle $.

For a state $\left| n \right\rangle $, what is expected value of $\left\langle x^2 \right\rangle $ and $\left\langle p^2 \right\rangle $? 
\begin{align*}
	\left\langle x^2 \right\rangle = \left\langle n \right| x^2 \left| n \right\rangle 
	&= \frac{\hbar}{2 m \omega} \left\langle n \right| (2N+1) \left| n \right\rangle  \\
	&= \frac{\hbar }{2 m \omega} (2 n + 1)
.\end{align*}
\begin{align*}
	\left\langle p^2 \right\rangle = \left\langle n \right| p^2 \left| n \right\rangle 
	&= - \frac{\hbar m \omega}{2} \left\langle n \right| \left(- a + a^{\dag}\right)^2 \left| n \right\rangle \\
	&= - \frac{\hbar m \omega}{2}\left\langle n \right|  \left(a^2 + \left(a^{\dag}\right)^2 - 2N - 1\right)\left| n \right\rangle  \\
	&= - \frac{\hbar m \omega}{2} (- 2 n - 1) \\
	&= \frac{\hbar m \omega}{2} (2 n + 1) 
.\end{align*}

Combining those, we have
\[
	\left\langle x^2 \right\rangle \left\langle p^2 \right\rangle 
	= \left( n + \frac{1}{2} \right) ^2 \hbar^2
.\] 
This implies the desired identity.
\end{proof}


